\documentclass{situationist-preprint}

\title{The Imitation Game Revisited}
\msid{MND-08-b}
\msfamily{MND-08}
\msdiscipline{Philosophy of Mind}
\mstype{Companion manuscript}
\mskeywords{Turing test; imitation game; functionalism; Blockhead; evidence for thought}

\begin{document}
\maketitle

\begin{abstract}
Turing proposed to replace the question whether machines can think by the question whether a machine could succeed in an imitation game. MND-08 argues that machines can think because thinking is a functional property, and accepts that behavior is not sufficient for thought. We ask what role remains for the imitation game. We argue that it is not a definition or a criterion of thinking but a test that provides defeasible evidence of the internal organization in which thinking consists. The evidence is strong when the test is open-ended, because the only feasible way to pass an open-ended test is to have such an organization. The Blockhead case shows that the evidence can be defeated, not that it is weak.
\end{abstract}

\section{Three readings of the game}

\citet{Turing1950} described a game in which an interrogator, communicating by text, tries to tell a machine from a human. He proposed that the question whether a machine could do well in the game replace the question whether machines can think. The proposal has been read in three ways. On the \emph{definitional} reading, to think is to be able to pass the test. On the \emph{criterial} reading, passing is sufficient for thinking, though thinking is not defined by it. On the \emph{evidential} reading, passing is evidence of thinking, defeasible by information about how the performance is produced \citep{Moor1976}.

MND-08-a rejects the first two. Thinking is a functional property (Principle~P1), had in virtue of internal states that stand in certain causal relations to inputs, outputs and one another (Definition~2.1), and a system that merely produces a thinker's outputs from a table does not have those states \citep{Block1981}. A definition or sufficient condition stated in terms of outputs alone is therefore false. The evidential reading remains.

\section{The evidential reading defended}

\begin{proposition}\label{p:evidence}
Success in an open-ended imitation game is strong evidence that the system has the functional organization of a thinker, defeasible by information about how the success is produced.
\end{proposition}

\begin{proof}
An open-ended game permits any question, on any topic, in any order, including questions that refer to earlier answers, ask for novel combinations, or require revising an earlier statement in the light of a new one. Success requires producing, for each such input, an output that a thinker might produce. There are two ways to do this. One is to have internal states that play the roles of beliefs, inferences and judgments, combining with one another to yield novel outputs, as Definition~2.1 of MND-08-a requires. The other is to store a sensible response for each possible conversation, as Block's table does. For a conversation of any interesting length the number of possible conversations exceeds what any physical system could store, so the second way is not available to a physically realized system. Success therefore makes the first way probable. The inference is defeated if we learn that the system works in some other way, for example that its answers are drawn from a stored list for a game short enough to tabulate.
\end{proof}

The proposition explains why Turing's test has seemed both compelling and inadequate. It is compelling because the open-endedness of conversation makes the internal organization of a thinker the best explanation of success. It is inadequate as a definition because the best explanation of a performance is not what the performance is. Blockhead is consistent with Proposition~\ref{p:evidence}: it describes a case in which the evidence is defeated by information about the mechanism.

\begin{objection}
Systems might succeed by a mechanism that is neither a table nor the organization of a thinker: a mechanism that reproduces the statistical shape of human conversation without states that play the roles of beliefs. Then success is not evidence of thinking at all.
\end{objection}

\begin{reply}
The objection describes a defeater, not a refutation. Whether a given mechanism has states that play the roles of beliefs and inferences is a question about its organization, to be settled by examining it, as Definition~2.1 of MND-08-a requires. If examination shows that it does, the system thinks; if it shows that it does not, the evidence from the game is defeated. Proposition~\ref{p:evidence} claims only that success raises the probability of the first answer when nothing else is known, and that is what the open-endedness of the test secures. As MND-08-a stresses, the conclusion that machines can think does not say which systems do.
\end{reply}

\begin{objection}
Human interrogators are easily misled, and a game judged by them measures the interrogator's credulity rather than the machine's capacities.
\end{objection}

\begin{reply}
The weakness lies in particular implementations of the game, not in its design. A game with expert interrogators, unlimited duration and the freedom to probe for the subcognitive associations that human thinkers share \citep{French1990} is harder to pass by shallow means, and success in it is correspondingly stronger evidence. The evidential reading accommodates differences in the quality of tests, as the other readings cannot.
\end{reply}

\section{Conclusion}

The imitation game is neither a definition nor a criterion of thinking. It is a test whose success, in its open-ended form, is strong but defeasible evidence of the functional organization in which thinking consists. Read this way, Turing's proposal is consistent with the conclusion of MND-08 and with the lesson of Blockhead: machines can think, and whether a given machine thinks is a question about its organization, on which its performance is evidence.

\begin{thebibliography}{9}\small
\bibitem[Block(1981)]{Block1981} Block, N. (1981). Psychologism and behaviorism. \emph{The Philosophical Review}, 90(1), 5--43.
\bibitem[French(1990)]{French1990} French, R.~M. (1990). Subcognition and the limits of the Turing test. \emph{Mind}, 99(393), 53--65.
\bibitem[Moor(1976)]{Moor1976} Moor, J.~H. (1976). An analysis of the Turing test. \emph{Philosophical Studies}, 30(4), 249--257.
\bibitem[Turing(1950)]{Turing1950} Turing, A.~M. (1950). Computing machinery and intelligence. \emph{Mind}, 59(236), 433--460.
\end{thebibliography}

\end{document}
